\section{Threats to Validity}\label{sec:threats}

\paragraph*{Construct validity}
The dependence set is measured with \numInfluencePerturbations{} random perturbations after every operation at
\numInfluenceInputs{} inputs, so it can miss true dependences but never adds false ones. The QDG's recall of
\numRqOneLcRecall{} is therefore a measured result that Proposition~\ref{prop:complete} guarantees in full, while the
distance between the QDG slice and the dependence set may be smaller than reported, since true dependences the
sampling missed would enlarge the set. That the dependence set is itself a sound slice was checked at
$\ket{0\dots0}$ only, which is why the paper does not call it a slice for every input. Wrong slices are counted at
$\ket{0\dots0}$ as well, a lower bound for every rule but the QDG, and RQ2 rechecks the slices under random inputs.
A slice preserves its criterion when the squared fidelity is at least $1 - 10^{-6}$, far above the oracle's error of
\numOracleErr{}. Each criterion of the study is a single qubit. The QDG slice of a set is the union of the slices of
its members by Proposition~\ref{prop:joint}, but this does not hold for the pruning rules, whose joint criteria were
explored only for state-based pruning. Fault containment measures where a developer would look, not whether a
developer finds the fault, and no developer took part in the study.

The QPDG is reimplemented from prose, which leaves room for readings its authors did not intend. Every reading was
logged before any run, and the flow-insensitive, flow-sensitive and read-only variants bracket the plausible ones.
Under the readings in which a gate writes its control, the QPDG contains the QDG slice, so its excess cannot come from
a reading that drops dependences. Commutation pruning is measured through Qiskit's pass, which applies to
\numLcqApplicable{} of the \numNPairs{} pairs, so its results describe circuits without mid-circuit measurements.

\paragraph*{Internal validity}
The slicers, the oracle and the dependence measurement could contain errors. The oracle agrees with Qiskit's
simulation to within \numOracleErr{}, a catalog of \numCatEntries{} hand-derived cases checks the ground truth and every
slicer before the runs, and the nesting of the dependence set within the QDG slice and of that within the QPDG slice,
which holds on all \numNPairs{} pairs, would expose an error that drops or adds operations in either. The soundness proof is mechanized for an
abstract model of local operations, and that quantum operations satisfy the model is argued on paper. The protocol was
written after an earlier run, so the outcomes of H1 and H2 were known, and both hold by construction. Every later
change is a dated deviation, and every result outside the five hypotheses, among them the classical PDG's share of
wrong slices and the nesting of the rules, is reported as exploratory. The language models were queried once per
item in the primary run. The Claude models accept only their default sampling, so their answers can differ between
calls, and the most expensive model was dropped before any scored call.

\paragraph*{External validity}
The circuits with an exact ground truth have at most \numSizeAMaxQubits{} qubits in corpus~A and 20 in corpus~L, and
only \numNDynamicPairs{} pairs come from dynamic circuits, so rules whose errors grow with width or with feedback may
fare worse on larger circuits than reported. The proofs hold for every circuit, but the measured distances between the
rules may change. Benchmark and generated circuits are not programs written by developers, and the \numLinesUnits{}
Bugs4Q units are the only real faults. Seeded faults follow published mutation operators and may not represent real
faults in their mix. The gate check covers Qiskit's standard two-qubit gates at one angle each, and the noise analysis
covers depolarizing noise only. Loops that run on the device are outside the model of Section~\ref{sec:core}, and
circuits that use them are not studied.
